%%%%%%%%%%%%%%%%%%%%%%%%%%%%%%%%%%%%%%%%%%%%%%%%%%%%%%%%%%%%%%%%%%%%%%%%%%%%%%%
%% assumption_contract.tex  -- the A1-A6 contract and the three coverage statements.
%%
%% PLACEMENT: immediately BEFORE the certificate theorem.
%%   kbound_short_body.tex : above % Finite-sample certificate: split out of theory_core_main.tex during the 2026-07 restructure.
% Consumed by the "what remains" method section (sec:method), not by the impossibility section.
\begin{proposition}[Coverage-to-action implication]\label{thm:certificate}\label{thm:cert}
Fix the scalar benefit target and its evaluation unit before calibration. Let $B$ and $\widehat B$
be finite real random variables, and let $\varepsilon\in[0,\infty]$ be a measurable random radius
on a common probability space containing all calibration, benefit-model fitting,
candidate-construction, and evaluation-unit randomness. The prediction $\widehat B$ and radius
$\varepsilon$ must be available without reading the new unit's evaluation labels. For
$\alpha\in(0,1)$, suppose the interval satisfies marginal coverage
\[
\Pr\!\bigl[\,|\widehat B-B|\le\varepsilon\,\bigr]\ \ge\ 1-\alpha .
\]
Define
\[
g=\begin{cases}
\adapt & \widehat B-\varepsilon>0,\\
\freeze & \widehat B+\varepsilon<0,\\
\abstain & \text{otherwise}.
\end{cases}
\]
Then, marginally over that common probability space,
\begin{align*}
\Pr\bigl[g=\adapt,B\le0\bigr]&\le\alpha,\\
\Pr\bigl[g=\freeze,B\ge0\bigr]&\le\alpha.
\end{align*}
\end{proposition}
\begin{proof}
On the event $\{|\widehat B-B|\le\varepsilon\}$, the condition
$\widehat B-\varepsilon>0$ implies $B>0$, so a false adapt requires the complement of
this event, which has probability at most $\alpha$. The freeze case is symmetric.
\end{proof}

\begin{remark}[Unit scaling: empirical radius versus theoretical budget]\label{rem:unit-scaling}
The empirical radius $\varepsilon$ and theoretical residual bound $\beta$ live on fundamentally
different scales. The population identity is $\Delta=2d(M+\gamma)$ where $d=P_T(D)>0$ is the
disagreement probability and $2d$ is twice that probability. Thus,
$|\Delta - 2dM| = 2d|\gamma| \le 2d\beta$. The unclipped population envelope centered at $2dM$
has radius $2d\beta$ on the benefit scale, not $\beta$; the exact population benefit set
$\Delta(\mathcal C_\beta)$ in Theorem~\ref{thm:frontier} includes clipping to $[-d, d]$:
$\Delta(\mathcal C_\beta) = 2d[\max(-\tfrac12, M-\beta), \min(\tfrac12, M+\beta)]$.
Equating $\varepsilon = \beta$ is therefore a conceptual category error: it conflates a conditional
score residual on the disagreement region with an empirical cell-benefit prediction error.
While numerical equality could occur by coincidence, their mathematical definitions remain distinct.
\end{remark}

\begin{proposition}[Cell-to-population transfer via compound bound]\label{thm:population-transfer}
Let $\Delta^{\mathrm{cell}}=\frac{1}{m}\sum_{j=1}^m W_j$ be the measured cell benefit on $m$ fresh,
conditionally i.i.d.\ evaluation draws, where
$W_j=\mathbf 1\{f_a(X_j)=Y_j\}-\mathbf 1\{f_0(X_j)=Y_j\}\in[-1,1]$, with conditional mean $\Delta$.
Suppose an independently fitted benefit predictor $h$ and split-conformal radius $\varepsilon$
achieve marginal cell coverage $\Pr(|\widehat\Delta-\Delta^{\mathrm{cell}}|\le\varepsilon)\ge 1-\alpha_{\mathrm{cell}}$
over exchangeable deployment episodes. For a sampling budget $\delta\in(0,1)$, define
$b(m,\delta)=\min\{2, \sqrt{2\log(2/\delta)/m}\}$. Then the compound interval
$I_{\mathrm{population}}=[\widehat\Delta-(\varepsilon+b),\; \widehat\Delta+(\varepsilon+b)]$ satisfies
\[
\Pr\bigl(\Delta\in I_{\mathrm{population}}\bigr)\ge 1-\alpha_{\mathrm{cell}}-\delta.
\]
The corresponding population decision rule $g_{\mathrm{pop}}$ using radius $\varepsilon+b$ satisfies
\begin{align*}
\Pr\bigl(\{g_{\mathrm{pop}}=\adapt, \Delta\le0\} &\cup {} \\
\{g_{\mathrm{pop}}=\freeze, \Delta\ge0\}\bigr) &\le \alpha_{\mathrm{cell}}+\delta.
\end{align*}
\end{proposition}
\begin{proof}
By Hoeffding's inequality for bounded variables in $[-1,1]$ of range 2,
$\Pr(|\Delta^{\mathrm{cell}}-\Delta|\le b)\ge 1-\delta$. On the intersection of cell coverage
and concentration, the triangle inequality yields $|\widehat\Delta-\Delta|\le |\widehat\Delta-\Delta^{\mathrm{cell}}|+|\Delta^{\mathrm{cell}}-\Delta|\le\varepsilon+b$.
A union bound proves coverage at level $1-\alpha_{\mathrm{cell}}-\delta$. Event containment
(as in Proposition~\ref{thm:certificate}) immediately bounds the probability of an incorrect strict
direction by $\alpha_{\mathrm{cell}}+\delta$.
\end{proof}

\begin{remark}[Direct empirical sampling of the population frontier]\label{rem:direct-sampled-frontier}
An alternative route directly operationalizes the population frontier $|M|>\beta$. Given fresh
unlabeled target inputs, let $n_D>0$ fall in $D$. The empirical disagreement margin
$\widehat M=\frac{1}{n_D}\sum_{X_i\in D}s(X_i)-\frac{1}{2}$ concentrates with radius
$r_M=\sqrt{\log(2/\alpha)/(2n_D)}$. Under the externally declared bound $|\gamma|\le\beta$, the rule
$\widehat M-r_M>\beta\Rightarrow\adapt$, $\widehat M+r_M<-\beta\Rightarrow\freeze$, and $\abstain$
otherwise, guarantees directional error at most $\alpha$. This couples sampling uncertainty $r_M$
additively with structural uncertainty $\beta$.
\end{remark}

\begin{remark}[Marginal vs.\ conditional error]\label{rem:fa-marginal}
Proposition~\ref{thm:certificate} bounds the \emph{marginal, unconditional} false-commit
probabilities for the declared target $B$. Its false-adapt rate is
$\mathrm{FA}_{\mathrm u}=\Pr[g=\adapt,B\le0]$. It does \emph{not} bound
$\mathrm{FA}_{\mathrm c}=\Pr[B\le0\mid g=\adapt]$, defined when
$\Pr[g=\adapt]>0$ and reported descriptively for $B=\Delta^{\mathrm{cell}}$.
Selection-conditional coverage requires a separate construction, such as that of
Jin and Ren~\cite{jin2025selection}; conditional guarantees over specified shift classes
also require their own calibration procedure~\cite{gibbs2025conditional}.
Neither is supplied by the marginal premise above. Simultaneous protection across candidates
and time-uniform protection likewise require additional guarantees.
Risk alignment (Definition~\ref{def:risk-align}) concerns population sign identifiability;
coverage for $B$ alone yields the \emph{interval} implication above.
\end{remark}

\begin{remark}[Fallback when assumptions are unsupported]\label{cor:abstain-valid}
If interval coverage for the declared target and evaluation unit---or the assumptions supporting
it---cannot be justified, Proposition~\ref{thm:certificate} supplies no level-$\alpha$ protection
for either strict action. A separately declared operational fallback may retain the frozen model
and record abstention. Serving the frozen model without a valid negative interval is not a
certified $\freeze$ decision. This policy recommendation does not establish that the retained
model is safe or optimal. If \emph{risk alignment} (Definition~\ref{def:risk-align}) cannot
be justified for an intended population sign claim, that interpretation remains unsupported.
The interval implication of Proposition~\ref{thm:certificate} still holds when coverage for the
declared $B$ is established (cf.\ Theorem~\ref{lem:nonid}).
\end{remark}
 (sec:method)
%%   kbound.tex            : above % ============================================================================
% FULL THEORY STACK — long paper (June 2026, reviewer-grade writing).
% Setup: paper/sections/theory_setup.tex (included from kbound.tex §\ref{sec:def}).
% ============================================================================

The body below expands the core spine into five load-bearing pillars.
\emph{Reproducibility:} numeric checks cited as ``validated'' are in the public
\texttt{theory\_v2/} scripts; they audit the algebra, not replace the written proofs.

\begin{lemma}[Disagreement-region reduction]\label{lem:reduction}\label{thm:disagree}
Under Assumption~\ref{ass:deploy}, with $\bar a:=\Pr_T(f_a{=}Y\mid D)$,
\[
\Delta = 2\,\mu_T(D)\bigl(\bar a-\tfrac12\bigr),
\qquad
\sign\Delta = \sign(\bar a-\tfrac12) = \sign(M+\gamma).
\]
For $K$-class $0/1$ loss, $\Delta=\mu_T(D)(p_a-p_0)$ on $D$; for squared loss the sign reduces
to one moment on $D$ (manuscript App.~C).
\end{lemma}
\begin{proof}
On $D$ exactly one predictor is correct under $0/1$ loss, giving
$\Delta=2\mu_T(D)(\bar a-\tfrac12)$. Linear expectation yields $\bar a-\tfrac12=M+\gamma$.
Multiclass and regression extensions are algebraic identities on $D$.
\end{proof}

\begin{lemma}[Plug-in regret identity]\label{lem:gate}\label{thm:gate}
For any estimator $\hat\Delta$ and gate $\hat g=\mathbf 1[\hat\Delta>0]$,
\[
R(\hat g)-R(g^\star)=\E\bigl[|\Delta(Z)|\,\mathbf 1\{\sign\hat\Delta\neq\sign\Delta\}\bigr],
\qquad g^\star=\mathbf 1[\Delta>0].
\]
On $\{|\Delta|<\varepsilon\}$, committal regret is at most $\varepsilon$.
\end{lemma}
\begin{proof}
Wrong-sign commits pay $|\Delta|$ pointwise; averaging gives the identity. On the low-margin
band, $|\Delta|\le\varepsilon$ bounds regret.
\end{proof}

\begin{proposition}[Finite-sample minimax regret floor]\label{prop:lecam-finite}\label{thm:imp-quant}
Let $\{P_\theta\}_{\theta\in\{-1,+1\}}$ be two target worlds with
$|\Delta(P_\theta)|\equiv\Lambda>0$, $\sign\Delta(P_\theta)=\theta$, and let $Q^{(n)}_\theta$
be the law of any label-free evidence from $n$ unlabeled observations. Then for every
label-free committal gate $\hat g$ and every finite $n$,
\[
\inf_{\hat g}\ \max_{\theta=\pm1}\bigl[R_{P_\theta}(\hat g)-R_{P_\theta}(g^\star)\bigr]
\ \ge\ \tfrac{\Lambda}{2}\Bigl(1-\TV\bigl(Q^{(n)}_{-1},Q^{(n)}_{+1}\bigr)\Bigr)
\ \ge\ \tfrac{\Lambda}{4}\,e^{-\mathrm{KL}(Q^{(n)}_{-1}\|Q^{(n)}_{+1})}.
\]
A Gaussian two-point family with $d_n=c/\sqrt n$ realizes a floor constant in $n$
(\texttt{val\_thm2\_lecam\_finite\_n.py}).
\end{proposition}
\begin{proof}
Lemma~\ref{lem:gate} with constant $|\Delta|=\Lambda$ reduces regret to the mixed committal
error $M(\hat g)$. Le Cam's testing affinity gives $\inf_{\hat g}M(\hat g)=1-\TV$; Bretagnolle--Huber
yields the KL bound.
\end{proof}

\begin{theorem}[Matched evidence impossibility]\label{thm:imp}\label{lem:nonid}
Fix $\alpha\in(0,\tfrac12)$.
\begin{enumerate}
\item[\textnormal{(i)}] \textbf{Witness.} There exist $P_T^1,P_T^2$ with
$\Law(Z\mid P_T^1)=\Law(Z\mid P_T^2)$, $\sign\Delta_1=-\sign\Delta_2\neq0$
(e.g.\ $X\sim\mathcal N(0,1)$, $f_a=1-f_0$, opposite label rules on $D$).
\item[\textnormal{(ii)}] \textbf{Le Cam floor.} For committal rules,
$\inf_g M(g)=1-\TV(P^{1,(n)}_{T,Z},P^{2,(n)}_{T,Z})$, equal to $1$ in the witness for all $n$.
\item[\textnormal{(iii)}] \textbf{Forced abstention.} Any rule with false-adapt and false-freeze
$\le\alpha$ must abstain on matched evidence with probability $\ge 1-2\alpha$.
\end{enumerate}
\end{theorem}
\begin{proof}
(i)~Construct worlds with common $X$-marginal and $Z$ a function of $X$ only, but opposite
benefit signs on $D$. (ii)~A committal rule is a test between the two evidence laws; apply Le Cam.
(iii)~Matched laws force identical action probabilities; each commit is wrong in one world.
\end{proof}

\begin{corollary}[Abstention is mandatory, not conservative]\label{cor:forced-abstain}
In the $|M|\le\beta$ wedge of Theorem~\ref{thm:frontier}, Theorem~\ref{thm:imp}(iii) applies:
any label-free rule with false-adapt and false-freeze $\le\alpha$ must abstain with probability at
least $1-2\alpha$ on matched evidence.
\end{corollary}

\begin{theorem}[Finite-sample certificate]\label{thm:cert}\label{thm:certificate}
Under Assumption~\ref{ass:deploy}, if calibration and deployment are exchangeable and
$\Pr[|\widehat\Delta-\Delta|\le\varepsilon]\ge1-\alpha$, then the three-way rule
(\textsc{adapt} if $\widehat\Delta-\varepsilon>0$, \textsc{freeze} if
$\widehat\Delta+\varepsilon<0$, else \textsc{abstain}) has marginal false-adapt and
false-freeze probabilities each at most $\alpha$.
\end{theorem}
\begin{proof}
On the coverage event $\{|\widehat\Delta-\Delta|\le\varepsilon\}$, the condition
$\widehat\Delta-\varepsilon>0$ implies $\Delta>0$; a false adapt requires the complement,
which has probability at most $\alpha$. The freeze case is symmetric.
\end{proof}

\begin{proposition}[Certificate sample complexity]\label{prop:cert-sample}
If $\widehat\Delta$ is an empirical-Bernstein mean of $n$ bounded paired benefits with variance
proxy $\sigma^2$ and range $b-a$, then one-sided certification ($\varepsilon<|\Delta|$) requires
$n=\tilde O(\sigma^2/\Delta^2)$; two-sided sign certification ($\varepsilon<|\Delta|/2$) requires
$n=\tilde O(\sigma^2/\Delta^2)$ with a larger constant. These rates are minimax order-optimal
(appendix; \texttt{val\_minimax\_optimality.py}).
\end{proposition}
\begin{proof}
Invert the empirical-Bernstein radius at level $\alpha$; the $8\sigma^2/\Delta^2$ constant
comes from requiring $\varepsilon<|\Delta|/2$.
\end{proof}

\begin{remark}[Sequential and multicandidate extensions]\label{rem:cert-ext}
Theorem~\ref{thm:cert} extends to anytime-valid streaming (Ville + betting supermartingale;
Theorem~\ref{thm:anytime}) and family-wise multicandidate routing (Bonferroni;
Theorem~\ref{thm:multicand}); proofs in Appendix~\ref{app:extensions}.
\end{remark}

\begin{theorem}[Exact benefit-sign frontier]\label{thm:frontier}\label{thm:headline}
Under Assumption~\ref{ass:deploy}, fix observables determining $M$ and budget $\beta\ge0$.
\begin{enumerate}
\item[\textnormal{(i)}] \textbf{Reduction.} $\sign\Delta=\sign(M+\gamma)$ for every $P_T$
(Lemma~\ref{lem:reduction}); $\gamma$ is one-dimensional given the observables.
\item[\textnormal{(ii)}] \textbf{Iff identifiability.} Over $\mathcal{C}_\beta$,
$\sign\Delta$ is a function of $\Law(Z)$ iff $|M|>\beta$; then $\sign\Delta=\sign M$ and the
minimax committal error is $0$. If $|M|<\beta$ the minimax committal error is $\tfrac12$ (opposite
nonzero signs share the evidence law); at $|M|=\beta$ the admissible drift $\gamma=-\sign(M)\beta$
makes $\Delta=0$, so strict-sign certification still fails and the maximal sound rule abstains on the
closed region $|M|\le\beta$.
\item[\textnormal{(iii)}] \textbf{Necessity geometry.} Any identifying class lies on one side of
the hyperplane $\gamma=-M$; every identifying assumption is a one-sided constraint on $\gamma$.
\end{enumerate}
The three-way rule $M>\beta\Rightarrow\textsc{adapt}$, $M<-\beta\Rightarrow\textsc{freeze}$,
else \textsc{abstain}, is the unique maximal sound certificate on $\mathcal{C}_\beta$.
\end{theorem}
\begin{proof}
Same as Theorem~\ref{thm:headline}: sufficiency from $|M|>\beta$ and $|\gamma|\le\beta$;
necessity from opposite drifts when $|M|\le\beta$ plus Theorem~\ref{thm:imp}; (iii) is
constancy of $\sign(M+\gamma)$ on an identifying class.
\end{proof}

\begin{remark}[$\beta=0$ face: prior label-free heuristics]\label{rem:beta-zero}
ATC~\cite{garg2022atc}, DoC, GDE~\cite{jiang2022gde}, COT~\cite{lu2023cot},
AETTA~\cite{lee2024aetta}, and agreement-on-the-line~\cite{baek2022aol,kim2024taol} threshold
observable surrogates of $M$ and commit to $\sign M$. They are sound on $\mathcal{C}_0$
($\gamma=0$) and share the failure mode $\gamma\neq0$ (Theorem~\ref{thm:frontier}(ii) at
$\beta=0$).
\end{remark}

\begin{corollary}[Knowable regimes]\label{cor:knowable-regimes}\label{thm:pos}
\textnormal{(a)} Under covariate shift with bounded density ratio, importance weighting
identifies $\sign\Delta$ from labeled source plus unlabeled target ($\beta=0$ on the weighted
score; premise not label-free testable). \textnormal{(b)} Conversely, no label-free certificate
identifies $\sign\Delta$ without a supplied bound on $|\gamma|$ (Theorem~\ref{thm:imp}).
\end{corollary}

\begin{theorem}[One-bit dichotomy]\label{thm:conj1-dichotomy}
Under the CEI panel hypothesis \textbf{H} on $D$ (conditionally independent correctness
indicators; manuscript for definitions):
\begin{enumerate}
\item[\textnormal{(i)}] \textbf{Rank-one law.} Under CEI alone, agreement statistics are
rank-one iff every per-class accuracy gap vanishes.
\item[\textnormal{(ii)}] \textbf{One global bit.} Under H, the evidence law fixes every
magnitude $|b_j|$ and product $b_ib_j$ but not the global orientation $b\mapsto -b$; the label
complement preserves $\Law(Z)$ while negating every benefit sign.
\item[\textnormal{(iii)}] \textbf{No evidence-definable bracketing.} A swap involution on $D$
fixes all label-free observable laws and negates $\Delta$; hence no testable assumption class
identifies $\sign\Delta$ without a bit-selector; the minimal untestable supplement is one bit.
\item[\textnormal{(iv)}] \textbf{Auditable budgets.} The bit-robust test of $|\gamma|\le\beta$
is level-$\alpha$ with power beyond $2r_n$; verification fails on an explicit blind set.
\end{enumerate}
Conjecture~\ref{conj:gen} (unconditional label-free bracketing) is resolved negatively by~(iii).
\end{theorem}
\begin{proof}[Proof sketch]
(i)~Algebraic expansion of agreements. (ii)--(iii)~Complement and swap involution on $D$;
observables span the even subalgebra while $\sign\Delta$ is odd. (iv)~Empirical-Bernstein
testing plus flip-minimization; full proofs in the supplementary appendix
(\texttt{val\_conj1\_*.py}).
\end{proof}

\begin{conjecture}[Label-free bracketing --- \textbf{resolved}]\label{conj:gen}
Unconditional label-free bracketing of the ordinal comparison on $D$ does not exist
(Theorem~\ref{thm:conj1-dichotomy}(iii)). The minimal supplement is one declared bit; weakest
one-bit classes are characterized in Appendix~\ref{app:weakest}
(Theorems~\ref{thm:cmono-weakest},~\ref{thm:uncond-weakest}).
\end{conjecture}

\begin{theorem}[Evidence-channel rate]\label{thm:ev-rate}
Under audited H with margin lower bounds, $m$ unlabeled samples on $D$ estimate agreement
statistics at rate $m^{-1/2}$ with constant $\Theta(1/(b_{\min}c_{\min}^2))$; labeled paired
benefits share the same exponent with efficiency ratio $\Theta(1/c_{\min})$. Labels buy constants
and the missing bit, not the rate (\texttt{val\_ev\_rate.py}).
\end{theorem}
\begin{proof}[Proof sketch]
Hoeffding on agreements, product-ratio perturbation, and a tensorized $\chi^2$/KL bound between
pattern laws differing in one coordinate; Bretagnolle--Huber finishes.
\end{proof}

\smallskip
\noindent\emph{Legacy cross-references.} Lemma~\ref{lem:reduction} is the disagreement-region
characterization formerly labeled \texttt{thm:disagree}; Lemma~\ref{lem:gate} is
\texttt{thm:gate}; Proposition~\ref{prop:lecam-finite} is \texttt{thm:imp-quant};
Corollary~\ref{cor:knowable-regimes} is \texttt{thm:pos}.
, or in the
%%                           Method section above the \paragraph{Assumptions.} block
%%
%% PORTABILITY: uses only amsmath, booktabs and the `definition' / `remark' theorem
%% environments, all present in kbound_tmlr.tex, kbound_short.tex and kbound.tex.
%% It deliberately declares NO algorithm float: kbound.tex loads algorithm2e while the
%% short drivers load algorithm+algpseudocode, and the two cannot both be satisfied by
%% one source file.
%%
%% DEFINES: sec:three-coverages, def:cov-theoretical, def:cov-observed,
%%          def:cov-diagnostic, def:inference-unit, ass:contract,
%%          rem:contract-not, rem:interval-carries, tab:fallback-ladder
%%%%%%%%%%%%%%%%%%%%%%%%%%%%%%%%%%%%%%%%%%%%%%%%%%%%%%%%%%%%%%%%%%%%%%%%%%%%%%%

\subsection{Three coverage statements, kept separate}\label{sec:three-coverages}

Much of the difficulty in reading an adaptation certificate comes from one word doing
three jobs. We fix the terminology here and hold to it for the rest of the paper. In
what follows $[L,U]=[\widehat\Delta-\varepsilon,\ \widehat\Delta+\varepsilon]$ is the
benefit interval of the certificate rule, so that the coverage premise
$\Pr[\,|\widehat\Delta-\Delta|\le\varepsilon\,]\ge1-\alpha$ of
Theorem~\ref{thm:certificate} and the interval statement below are the same hypothesis
written two ways.

\begin{definition}[Theoretical coverage]\label{def:cov-theoretical}
The interval has \emph{theoretical coverage} at level $1-\alpha$ if
\begin{equation}\label{eq:cov-theoretical}
  \Pr\bigl\{\Delta \in [L,U]\bigr\} \;\ge\; 1-\alpha ,
\end{equation}
the probability taken over the common space of calibration, benefit-model fitting and
deployment-unit randomness. This is a property of a distribution, not of a run. We
assert it only where (A1)--(A6) below are argued to hold, and never on the basis of an
observed hit rate.
\end{definition}

\begin{definition}[Observed empirical coverage]\label{def:cov-observed}
For a completed, labelled track with $n$ independent inference units
(Definition~\ref{def:inference-unit}), the \emph{observed empirical coverage} is
\begin{equation}\label{eq:cov-observed}
  \widehat{\mathrm{Cov}}
  \;=\;
  \frac{1}{n}\sum_{i=1}^{n}\mathbf{1}\bigl\{\Delta_i \in [L_i,U_i]\bigr\},
\end{equation}
reported with $n$, the resampling unit, and an interval computed at that unit. It
estimates a hit rate on the conditions evaluated, and is not evidence for
\eqref{eq:cov-theoretical} on any other condition.
\end{definition}

\begin{remark}[When \emph{not} to attach an interval]\label{rem:no-interval}
Definition~\ref{def:cov-observed} applies to a hit rate that is genuinely random. It does
\emph{not} apply to a realized calibration-set coverage figure produced by in-sample rank
calibration, which is a deterministic function of $n$ and carries no sampling variability
at all --- attaching a binomial or bootstrap interval to such a figure gives it a
frequentist meaning it does not have. \S\ref{sec:limits} identifies the figures in this
paper that are arithmetic in this sense, and they are reported without intervals. The
discipline this section imposes is that every coverage number states which of the three
kinds it is; for the arithmetic ones, the honest annotation is ``deterministic in $n$'',
not a wider interval.
\end{remark}

\begin{definition}[Deployment diagnostic]\label{def:cov-diagnostic}
A \emph{deployment diagnostic} is a label-free statistic whose large values are evidence
\emph{against} (A1)--(A6). Diagnostics are one-sided by construction: they can falsify
and cannot verify. Passing every diagnostic of \S\ref{sec:gate} does not establish
\eqref{eq:cov-theoretical}.
\end{definition}

\begin{definition}[Inference unit]\label{def:inference-unit}
The \emph{inference unit} is the level at which calibration draws are treated as
exchangeable and at which uncertainty is resampled: a domain, an episode, a
corruption$\times$severity cell, a backbone, or a seed, declared per track. It is the
individual example only when examples are independently drawn, which on the tracks of
\S\ref{sec:exp} they are not.
\end{definition}

\subsection{The assumption contract}\label{ass:contract}

The paper already records which condition buys what, and what to do when one
fails. What that table does not do is give the conditions names that the
rest of the paper, the released artifacts and a deployment can all refer to. We do that
now, because a condition nobody can cite is a condition nobody checks.

\begin{description}
  \item[\textnormal{(A1)} Calibration exchangeability.] Calibration units and the
    deployment unit are exchangeable at the declared inference unit, or an explicit
    shift correction supplies the same conclusion. (This is the episode-exchangeability condition E2, read at the unit level.)
  \item[\textnormal{(A2)} Risk alignment.] The evidence $Z$ is risk-aligned
    (Definition~\ref{def:risk-align}): it retains enough information about $\Delta$ for
    calibration residuals to transfer to the deployment condition.
  \item[\textnormal{(A3)} Coverage premise.] The interval satisfies
    \eqref{eq:cov-theoretical} on the target class.
  \item[\textnormal{(A4)} Protocol independence.] Target-test labels influence none of:
    feature selection, candidate selection, threshold selection, calibration, or
    stopping. (The retrospective outcome-logging condition E1, and the data-access
    discipline of the protocol tables.)
  \item[\textnormal{(A5)} Correct sampling unit.] Calibration and uncertainty
    calculations operate on independent domains, episodes, corruptions, backbones or
    seeds, not on correlated examples pooled as though independent.
  \item[\textnormal{(A6)} Fixed decision rule.] Estimator, interval construction,
    $\alpha$ and decision rule are fixed before the promoted target conditions are
    evaluated.
  \item[\textnormal{(A7)} Estimator stability.] Where the radius is calibrated
    leave-one-out rather than on a clean split, the deployed full-data fit and each
    leave-one-out fit differ by at most $\beta_{\mathrm{stab}}$ in prediction.
\end{description}

\noindent (A1)--(A4) and (A6) rename conditions the paper already carried:
the conditions previously listed as risk-aligned $Z$, coverage of
$\widehat\Delta\pm\varepsilon$, exchangeable or corrected residuals, and frozen
protocol become (A2), (A3), (A1) and (A4)/(A6) respectively, and the long manuscript's assumptions (i)--(iii) map to (A1), the
bounded-benefit regularity condition, and (A2). Bounded loss and benefit remains a
regularity condition needed for the concentration arguments; it is not part of the
deployment gate because it is a property of the loss we choose, not of the deployment we
meet. \textbf{(A5) is new in this revision}, and it is the one the empirical sections
most often failed to make explicit: a coverage figure computed over correlated cells was
reported without the denominator that makes it readable.

\begin{remark}[Why (A7) exists, and what it costs]\label{rem:a7}
The grids of \S\ref{sec:exp} calibrate $\varepsilon$ leave-one-out rather than on a
clean split, and no theorem makes that exact conformal: the plain jackknife has no
distribution-free finite-sample guarantee, and its finite-sample counterpart buys
$1-2\alpha$ rather than $1-\alpha$. What does hold is a deterministic transfer. If the
deployed fit differs from each leave-one-out fit by at most $\beta_{\mathrm{stab}}$,
then leave-one-out coverage at radius $\varepsilon$ gives full-fit coverage at radius
$\beta_{\mathrm{stab}}+\varepsilon$ (\path{formal/KBound/Stability.lean},
\texttt{stability\_transfers\_loo\_coverage}, machine-checked). So
$\beta_{\mathrm{stab}}$ is the price of leave-one-out calibration, and it is an
assumption about the estimator rather than something the calibration data can
establish. The contrapositive is the useful direction: a miss beyond
$\beta_{\mathrm{stab}}+\varepsilon$ \emph{proves} the fit was less stable than
$\beta_{\mathrm{stab}}$, so (A7) is falsifiable after the fact even though it is not
verifiable before.
\end{remark}

\begin{remark}[What the contract is not]\label{rem:contract-not}
Theorem~\ref{thm:certificate} is conditional on (A1)--(A6), and K-Bound does not infer,
test into, or prove any of them from unlabelled deployment data. (A1)--(A3) are not
identifiable label-free; this is the obstruction of \S\ref{sec:theory} applied to
calibration transfer rather than to the drift budget. What a deployment can do is detect
\emph{some} violations and refuse to certify when it does, which is \S\ref{sec:gate}.
\end{remark}

\begin{remark}[Where the guarantee lives]\label{rem:interval-carries}
The proof of Theorem~\ref{thm:certificate} uses no property of $\widehat\Delta$ beyond
its appearance in $L$ and $U$: a false adapt requires $\Delta\le0<L$, which is the
complement of the coverage event, and a false freeze requires $\Delta\ge0>U$. So the
guarantee is carried entirely by (A3); the fitted benefit estimator contributes decision
\emph{yield}, not safety. A better estimator narrows the interval and buys more committed
decisions at the same $\alpha$; a worthless one widens it until the system abstains. That
is the intended failure direction, and it is why (A3) --- not estimator quality --- is
the assumption a deployment has to defend.
\end{remark}

\subsection{Permitted behaviour by assumption state}\label{sec:fallback-ladder}

Where the earlier tables say what to do when a single condition fails,
Table~\ref{tab:fallback-ladder} says what a deployment is permitted to \emph{emit},
which is the operational form of the same idea and the object \S\ref{sec:gate}
implements.

\begin{table}[htbp]
\centering
\small
\caption{Fallback ladder. Permitted behaviour is a function of the assumption state, not
of the estimated benefit. ``Certification disabled'' means no certificate is emitted; the
system may still serve the frozen model, but no K-Bound claim attaches to that choice.}
\label{tab:fallback-ladder}
\begin{tabular}{@{}p{0.46\columnwidth}p{0.44\columnwidth}@{}}
\toprule
Assumption state & Permitted behaviour \\
\midrule
All required gates pass
  & Full \textsc{adapt}/\textsc{freeze}/\textsc{abstain} certificate \\
Support warning, no hard violation
  & \textsc{abstain} or \textsc{freeze} only \\
Insufficient independent calibration units
  & Diagnostic output only; no decision \\
Leakage, failed transfer, or unknown provenance
  & Certification disabled \\
\bottomrule
\end{tabular}
\end{table}

Read with Remark~\ref{rem:interval-carries}, this is the honest description of what
K-Bound offers: \emph{conditional insurance against detectable adaptation harm}, not an
unconditional safety guarantee.
